\section{The actual invariant kernel: degree filtration, quantitative minors, and its complete conductor graph}\label{the-actual-invariant-kernel-degree-filtration-quantitative-minors-and-its-complete-conductor-graph}

21 September 2026. Independent derivation on the original simple-quartet five-orbit domain. All matrices below use the original period, amplitude, root ordering and units. In particular the invariant image is the entire matrix \texttt{Z={[}M\_A,S\_k\ Z\_k{]}}, not just its conductor columns.

The proved advancement is a quantitative degree-filtration theorem and a reduction of the actual kernel determinant, with error \texttt{o(kq)}, to its specified image in the complete weighted conductor-relation quotient. The fixed low-degree graph component is controlled by an actual angle estimate from invariant jets. A rank estimate alone would not prove this. The growing relation operator remains present in the final receiver. Its replacement by an outer polynomial ideal is not asserted.

\subsection{1. Original data and the exact residual constraint}\label{original-data-and-the-exact-residual-constraint}

Retain \texttt{k\textgreater{}=9}, \texttt{k=1\ modulo\ 4}, \[
q=(k+1)^2,\quad k'=k-8,\quad q'=(k-7)^2,\quad
\Delta=q-q'=16k-48,\quad m=8k-16,\quad r_0=8k-32.
\tag{KF1}
\] The original roots and their lower counterparts are \[
\beta_{ab}=k/2+(2a-k)\delta+i(2b-k)\gamma,
\quad
\beta'_{ab}=(k-8)/2+(2a-k+8)\delta+i(2b-k+8)\gamma,
\] with \texttt{0\textless{}delta\textless{}1/2}, \texttt{gamma\textgreater{}2}. Let chi and chi' be their full monic root polynomials. The literal value matrix is V, and the full original kernel frame is \[
I_K=V^{-1}\pi^T,\quad
VK=\ker Z^T=\operatorname{im}\pi^T,
\quad Z=[M_A,S_kZ_k].
\tag{KF2}
\] Throughout, the four-return means exactly \nolinkurl{R f=f_{q-1}+f_q-f_{2q-1}-f_{2q}}. Transposes in this equality are complex-linear transposes. Metric adjoints will be written separately. The full conductor biform is \nolinkurl{A=sum_{r,s=0}^8 a_{rs}e_{rs}^{(8)}}. Define \[
b_{rs}=4+(2r-8)\delta+i(2s-8)\gamma,\quad
E_A(z)=\sum_{r,s}a_{rs}e^{b_{rs}z},\quad
\mu_j=E_A^{(j)}(0),\quad v=\operatorname{ord}_0E_A\le80.
\tag{KF3}
\] The constant v and the nonzero mu\_v are fixed by the original data. Put \[
g=\Delta-v,\qquad d=q'+v,\qquad
\mathcal L=\mathbb C[S]_{<v}\subset E=\mathbb C[S]/(\chi).
\] When v=0 this space is zero and every statement concerning its angles or jets below is empty.

The complete conductor operator is \[
\mathcal T_Ap(S')=\sum_{r,s}a_{rs}p(S'+b_{rs})
=\frac{\mu_v}{v!}\partial^v B(\partial)p(S'),\quad
B(z)=\frac{v!E_A(z)}{\mu_vz^v},\quad B(0)=1.
\tag{KF4}
\] On any fixed polynomial space this is a finite identity: only the derivatives up to the degree of p occur. It kills exactly the polynomials of degree below v and lowers every other degree by exactly v, with \[
[S'^{n-v}]\mathcal T_A S^n=\mu_v\binom nv\ne0.
\tag{KF5}
\] Root addition proves \nolinkurl{chi' divides T_A(chi h)} for every h. The conductor kernel is \[
V_A=\{[p]:\chi'\mid\mathcal T_Ap\},\quad \dim V_A=\Delta,
\qquad K\subset V_A.
\tag{KF6}
\] For a fixed representative of degree below q define \[
\eta(p)=\mathcal T_Ap/\chi'\in\mathbb C[S']_{<g},\quad p\in V_A.
\tag{KF7}
\] It is a polynomial by KF6. The quotient map \texttt{V\_A-\textgreater{}P\_\{\textless{}g\}} is onto and has kernel L: onto follows by successively solving KF5 for a preimage of \texttt{chi\textquotesingle{}\ eta}, whose degree is below q-v; the kernel follows from KF4. This is the cutoff-independent low-polynomial version of the CEP43 map, with its unscaled mu\_v factor retained.

Let R be the unique right inverse of T\_A with zero coefficients in degrees below v. An explicit construction is \[
R f=\frac{v!}{\mu_v}
 \left[B(\partial)^{-1}\partial^{-v}_0f
 -\operatorname{low}_{<v}(B(\partial)^{-1}\partial^{-v}_0f)\right].
\tag{KF8}
\] Here \nolinkurl{partial_0^{-v}S'^j=j!S^{j+v}/(j+v)!}, and \texttt{B(partial)\^{}\{-1\}} is its finite nilpotent inverse on the degree in question. Subtracting the low part changes no image under T\_A, so KF8 has the asserted properties. Thus every original representative in V\_A is uniquely \[
p=l+R(\chi'\eta),\qquad l\in\mathcal L,\quad\deg\eta<g.
\tag{KF9}
\] Let J\_R be the complete matrix \texttt{S\_kZ\_k} of the additional r\_0 invariant columns. The actual remaining constraint is exactly \[
\mathcal B_k(l,\eta):=J_R^TV\bigl(l+R(\chi'\eta)\bigr)=0.
\tag{KF10}
\] It has rank r\_0 on the Delta-dimensional space in KF9, since its kernel is exactly K of dimension m. The conductor columns were used to reach V\_A, and every other invariant column remains in KF10. There are no generic coefficients in this description.

\subsection{2. Exact degree filtration and high-degree projection minors}\label{exact-degree-filtration-and-high-degree-projection-minors}

For \texttt{v\textless{}=n\textless{}d=q\textquotesingle{}+v}, a polynomial p of degree at most n in V\_A has \texttt{deg\ T\_Ap\textless{}q\textquotesingle{}}. Divisibility by the monic polynomial chi' of degree q' therefore forces T\_Ap=0, so p belongs to L. Hence \[
V_A\cap\mathcal P_n=\mathcal L\quad(v-1\le n<d),\qquad
K\cap\mathcal P_n=K\cap\mathcal L\quad(v-1\le n<d).
\tag{KF11}
\] For \texttt{0\textless{}=t\textless{}g}, the leading coefficient in KF5 gives \[
\deg R(\chi'\eta)=q'+v+\deg\eta\quad(\eta\ne0).
\] Consequently the complete actual filtration is \[
\dim(K\cap\mathcal P_{d+t})
=v+t+1-\operatorname{rank}
 \mathcal B_k\big|_{\mathcal L\oplus\mathcal P'_t}.
\tag{KF12}
\] The initial value is \nolinkurl{s_k=dim(K intersect L)=v-rank B_k|_L}. All these ranks use the actual matrix KF10. At the last t=g-1 their value is m. This proves that every degree pivot outside the first v positions lies in the last g positions \texttt{d,...,q-1}. It does not assert that the last m positions alone give an invertible minor.

There is a quantitative coefficient comparison on this high block. For p in V\_A write \texttt{h\_j={[}S\^{}\{d+j\}{]}p}, \texttt{0\textless{}=j\textless{}g}, and let eta\_j be the coefficients in KF7. Let C be the high-coefficient multiplication matrix of chi': its entry \texttt{C\_\{ij\}} for i\textless=j is the coefficient of \texttt{S\textquotesingle{}\^{}\{q\textquotesingle{}-(j-i)\}} in chi', and it is zero for i\textgreater j. Let T have entries \[
T_{ij}=\binom{d+j}{v+j-i}\mu_{v+j-i}\quad(i\le j),\qquad
T_{ij}=0\quad(i>j).
\tag{KF13}
\] Comparison of coefficients in \texttt{T\_Ap=chi\textquotesingle{}\ eta} proves \[
C\eta=Th,\qquad \eta=D_kh,\quad D_k=C^{-1}T,
\qquad
\det D_k=\mu_v^g\prod_{j=0}^{g-1}\binom{d+j}{v}.
\tag{KF14}
\] C has diagonal one. This determinant is nonzero for the actual data, including v=0, and has not been inferred from genericity.

The following explicit bounds make its conditioning quantitative. Put \[
A_{\rm abs}=\sum|a_{rs}|,\quad B_0=\max(1,\max|b_{rs}|),
\quad C_v=A_{\rm abs}B_0^v/|\mu_v|,
\] \[
R_0=B_0^{-1}\min\{1,(v+1)/(2eC_v)\},\qquad
R'=\max(1,\max|\beta'_{ab}|).
\tag{KF15}
\] The exponential Taylor series gives \nolinkurl{|B(z)-1|<=C_v B_0|z| exp(B_0|z|)/(v+1)<=1/2} for \textbar z\textbar\textless=R\_0. Thus the coefficients c\_l of 1/B obey \nolinkurl{|c_l|<=2R_0^{-l}} by Cauchy's coefficient integral. The complete inverse of T is \[
(T^{-1})_{ij}=\frac{v!}{\mu_v}c_{j-i}
 \frac{(q'+j)!}{(d+i)!}\quad(i\le j).
\tag{KF16}
\] This follows by applying the integration and inverse derivative series in KF8 to a single high output monomial; lower terms do not enter the high coefficient block. Since the factorial ratio is at most q\^{}\{j-i\}, the following conservative Euclidean operator bounds hold: \[
\begin{split}
\|C\|&\le g(1+q'R')^{g-1},\\
\|C^{-1}\|&\le g(1+(q'+g)R')^{g-1},\\
\|T\|&\le gA_{\rm abs}(1+qB_0)^{v+g-1},\\
\|T^{-1}\|&\le\frac{2gv!}{|\mu_v|}
                 \max(1,q/R_0)^{g-1}.
\end{split}\tag{KF17}
\] For C, use the elementary symmetric coefficient bound \nolinkurl{binom(q',l)(R')^l}; for its inverse use the complete homogeneous coefficient bound \nolinkurl{binom(q'+l-1,l)(R')^l}. The matrix norm is at most g times its largest entry. The T bound uses the finite binomial sum and \nolinkurl{|mu_j|<=A_abs B_0^j}. These proofs retain every coefficient.

Let \nolinkurl{B_+=||C^{-1}||_bound ||T||_bound} and \nolinkurl{B_-=||T^{-1}||_bound ||C||_bound}, where the bounds mean exactly KF17. Then \[
\|D_k\|\le B_+,\quad\|D_k^{-1}\|\le B_-,\qquad
\log^+B_++\log^+B_-=O_{\rm actual}(k\log(k+1)).
\tag{KF18}
\] For any h-dimensional actual plane in V\_A/L, let H and Y be its high-coefficient and eta-coefficient matrices in the same frame. KF14 says Y=D\_kH. Cauchy--Binet and the operator bounds give \[
B_+^{-2h}\det(Y^*Y)
\le\sum_{|J|=h}|\det H_J|^2
\le B_-^{2h}\det(Y^*Y).
\tag{KF19}
\] In particular at least one actual high-degree projection minor satisfies \[
|\det H_J|^2\ge
\binom gh^{-1}B_+^{-2h}\det(Y^*Y)>0.
\tag{KF20}
\] The statement applies to the image of K after quotient by K intersect L, with \texttt{h=m-s\_k}. Its logarithmic coefficient distortion is \nolinkurl{O_actual(k^2 log(k+1))=o(kq)}. The Euclidean coefficient metrics used in KF17--KF20 have been specified explicitly. They are not identified with a Gamma minimum metric; the metric transport below is separate.

\subsection{3. The actual low-degree angle from the full invariant image}\label{the-actual-low-degree-angle-from-the-full-invariant-image}

The OCF data give an invertible matrix H\_mathsf and the exact substitution \nolinkurl{x=H_mathsf(t_1,t_2,t_3,t_4)^T}. Let rho\_1,\ldots,rho\_4 be the original one-factor roots in their prescribed order, and put \[
x_i(z)=\sum_{j=1}^4(H_{\rm mathsf})_{ij}e^{\rho_jz}.
\tag{KF21}
\] A degree-k invariant monomial F(x) restricts to the complete biform coefficient vector z\_F and hence satisfies \[
F(x(z))=\sum_{a,b}(z_F)_{ab}e^{\beta_{ab}z},\qquad
z_F^TVp=\left.p(\partial_z)F(x(z))\right|_{z=0}.
\tag{KF22}
\] The first equality is the original Segre multiplication rule; its second equality follows by differentiating each term. This is the exact link from invariant jets to the low polynomial directions.

Assume v\textgreater=1. Work in the finite jet algebra \texttt{J\_v=C{[}z{]}/(z\^{}v)}. Let J\_+ be the indices i with x\_i(0) nonzero, and J\_0 its complement. There is at least one index in J\_+ because H\_mathsf is invertible and the vector \texttt{(1,1,1,1)} is nonzero. Choose any such p.~Write \texttt{s=\textbar{}J\_+\textbar{}}, \texttt{z\_0=\textbar{}J\_0\textbar{}=4-s}.

For a fixed residue epsilon modulo five define the finite base set B\_epsilon as the exponent vectors r satisfying all of the following:

\begin{itemize}
\tightlist
\item
  \texttt{0\textless{}=r\_i\textless{}=4} for i in J\_+;
\item
  \texttt{0\textless{}=r\_i\textless{}=v-1} for i in J\_0;
\item
  \nolinkurl{sum_i i r_i=0 modulo 5} and \nolinkurl{sum_i r_i=epsilon modulo 5}.
\end{itemize}

The weights i are exactly the original C5 weights 1,2,3,4. If this base set is empty, every degree-k invariant jet is zero for k in this residue, so L is contained in K; no nonzero low angle is then needed. Otherwise let \[
D_\epsilon=\max_{r\in B_\epsilon}|r|,\quad
A_v=(s-1)(v-1),\quad d_\epsilon=D_\epsilon+5A_v.
\tag{KF23}
\] It obeys \nolinkurl{d_epsilon<=4s+(v-1)(4-s)+5(s-1)(v-1)<=15v+1}. Every base degree has the same residue as d\_epsilon.

For i in J\_+ other than p define the unit jet \texttt{R\_i=(x\_i/x\_p)\^{}5}. Its difference from its constant term is nilpotent of order at most v. Thus \nolinkurl{(R_i-R_i(0))^v=0}, and all powers of R\_i are linear combinations of \nolinkurl{1,R_i,...,R_i^{v-1}}. Products of these powers reduce to exponents \texttt{0\textless{}=a\_i\textless{}v} with total at most A\_v.

For every r in B\_epsilon and every such exponent list a, the following is an actual invariant monomial of degree d\_epsilon: \[
F_{r,a}(x)=x^r\prod_{i\in J_+\setminus\{p\}}x_i^{5a_i}
 x_p^{d_\epsilon-|r|-5\sum a_i}.
\tag{KF24}
\] Its last exponent is a nonnegative multiple of five by KF23. These give a fixed finite list of monomials, independent of growing k.

For every \texttt{k\textgreater{}=d\_epsilon} of that residue, the complete space of degree-k invariant jets is exactly \[
\boxed{\operatorname{span}j_{<v}\{F(x(z)):F\in\mathscr I_k\}
=j_{<v}(x_p^{k-d_\epsilon})
 \operatorname{span}_{r,a}j_{<v}F_{r,a}(x(z)).}
\tag{KF25}
\] To prove the nontrivial inclusion, first discard any monomial having an exponent at least v on a zero-constant x\_i: its jet vanishes. For every nonzero-constant variable divide its exponent by five, retaining its residue between zero and four. The resulting base lies in B\_epsilon. Factor the remaining powers as a common power of x\_p times powers of the unit ratios R\_i. Reduce each latter power by its nilpotent polynomial identity. KF24 is precisely the resulting finite list after extracting the common factor \nolinkurl{x_p^{k-d_epsilon}}. Conversely each monomial in KF24 multiplied by \nolinkurl{x_p^{k-d_epsilon}} is itself invariant and has degree k, so its jet belongs to the left side. This proves equality.

Write J\_epsilon for the v by J matrix whose columns are the Taylor coefficients of the fixed jets in KF24. Let t\_epsilon be its rank. If it is zero, the preceding empty-angle case again applies. Otherwise let \[
s_\epsilon=\sigma_{\min}^+(J_\epsilon)>0,\quad
a=x_p(0)\ne0,\quad
w_* =\sum_{j=1}^{v-1}\frac{|x_p^{(j)}(0)|}{j!|a|},
\quad n=k-d_\epsilon.
\tag{KF26}
\] All of these are fixed finite expressions in the original period and roots except n.~The positive constant s\_epsilon can equivalently be defined by the square root of the smallest positive eigenvalue of \nolinkurl{J_epsilon J_epsilon^*}; no numerical evaluation or generic lower bound is claimed for it. Its positivity follows from its actual nonzero rank.

Multiplication by the unit jet x\_p\^{}n is an invertible matrix U\_n.~The nilpotent binomial inverse gives the explicit bound \[
\|U_n^{-1}\|\le B_n^{\rm unit}:=
|a|^{-n}\sum_{j=0}^{v-1}\binom{n+j-1}{j}w_*^j\quad(n\ge1),
\qquad B_0^{\rm unit}=1.
\tag{KF27}
\] Indeed write the multiplication matrix of x\_p as \texttt{a(I+W\_0)}, with \texttt{W\_0\^{}v=0} and \texttt{\textbar{}\textbar{}W\_0\textbar{}\textbar{}\textless{}=w\_*}; expand \texttt{(I+W\_0)\^{}\{-n\}} exactly through degree v-1. In particular \nolinkurl{log^+B_n^unit=O_actual(k+log(k+1))}.

Use the actual invariant monomials \nolinkurl{F_{r,a}(x)x_p^{k-d_epsilon}} as rows of a coefficient matrix B\_k\^{}T. Put \[
H_* =\max\left(1,\max_i\sum_j|(H_{\rm mathsf})_{ij}|\right),
\quad R_k^S=\max(1,\max|\beta_{ab}|),\quad
L_k^{\rm val}=\sqrt{qv}(R_k^S)^{v-1}.
\tag{KF28}
\] The multinomial theorem and the complete Segre substitution give \nolinkurl{||B_k^T||<=sqrt(J)H_*^k}; there are no coefficients deleted in this estimate. The value map on coefficient vectors of polynomials of degree below v has norm at most L\_k\^{}val. By KF22 its low observation matrix is exactly \[
O_k=J_\epsilon^T U_n^T\operatorname{diag}(0!,1!,\ldots,(v-1)!).
\tag{KF29}
\] Ordinary transpose preserves singular values, and the smallest singular value of the factorial matrix is one. Thus \[
\sigma_{\min}^+(O_k)\ge s_\epsilon/B_n^{\rm unit}.
\tag{KF30}
\] For clarity, put \texttt{A=J\_epsilon\^{}T} and \nolinkurl{B=U_n^T diag(0!,..., (v-1)!)}. Then \nolinkurl{BB^* >= (B_n^unit)^{-2} I}, so \nolinkurl{O_k O_k^* >= (B_n^unit)^{-2} AA^*}. The two products have the same rank because B is invertible. The minimum characterization of their ordered eigenvalues proves KF30.

KF25 proves that the kernel of this selected low observation is exactly K intersect L, rather than the kernel of a smaller selection. For l in L, let distance be measured in Euclidean root-value coordinates. Then \[
\boxed{\operatorname{dist}(Vl,VK)
\ge\varepsilon_k\operatorname{dist}(Vl,V(K\cap\mathcal L)),}
\quad
\varepsilon_k=
\min\left(1,\frac{s_\epsilon}
 {B_n^{\rm unit}\sqrt JH_*^k L_k^{\rm val}}\right)>0.
\tag{KF31}
\] To prove it, every row of B\_k\^{}T vanishes on VK, so its value on Vl is at most \nolinkurl{sqrt(J)H_*^k dist(Vl,VK)}. On the other hand KF30 bounds it below by \nolinkurl{(s_epsilon/B_n^unit)} times coefficient distance to its kernel. That coefficient distance is at least the root-value distance to V(K intersect L) divided by L\_k\^{}val. Combining the two estimates gives KF31. Its logarithmic cost is \nolinkurl{log(1/epsilon_k)=O_actual(k+log k)}. For the finitely many k below d\_epsilon the same statement uses the finite matrix of all degree-k invariant monomials in place of KF29; its positive singular values give an exact finite epsilon\_k. The asymptotic statement needs only the proved stabilized range. When v=0 or the selected observation has rank zero, set epsilon\_k=1. In the latter case L is contained in K, and both distances in KF31 are zero. These conventions include every endpoint and residue class.

It also follows that \[
\dim(K\cap\mathcal L)=v-t_\epsilon
\quad(k\ge d_\epsilon,\ k=\epsilon\bmod5).
\tag{KF32}
\] This is an exact residue-class statement about the actual period.

\subsection{4. Why top-degree transversality cannot be assumed from the five-orbit condition}\label{why-top-degree-transversality-cannot-be-assumed-from-the-five-orbit-condition}

The stronger assertion \texttt{K\ intersect\ L=0} is not a consequence of the algebraic five-orbit hypothesis alone. An explicit auxiliary matrix is \[
H_{\rm ex}=\begin{pmatrix}
1&0&0&0\\-1&1&0&0\\-1&0&1&0\\-1&0&0&1
\end{pmatrix}.
\] It is invertible and sends \texttt{(1,1,1,1)\^{}T} to the first coordinate vector. Its transformed Segre equation is \[
Q=x_1x_4-x_1x_2-x_1x_3-x_2x_3.
\] It is an irreducible rank-four quadric. Under the diagonal C5 action, its nonzero terms have weights 0,3,4, so its five translates are pairwise nonassociate: an associate would have scalar one from the weight-zero coefficient and would require both \nolinkurl{zeta^{3a}=zeta^{4a}=1}. At the first coordinate vector every invariant homogeneous monomial of degree k vanishes unless k is divisible by five. Therefore the constant polynomial has zero pairing with the entire invariant image for all k not divisible by five, and belongs to that example's actual kernel. Its conductor symbol vanishes at zero as well, since each translated quadric vanishes at that point, so this direction lies in its space L.

This is a counterexample within the stated invertible-matrix five-orbit algebra, not an assertion that H\_ex is the programme's period matrix. For the original period KF25--KF32 give the exact deciding map and its quantitative nonzero constants. They do not replace it by the example.

\subsection{5. Uniform Gamma condition bounds in the original value coordinates}\label{uniform-gamma-condition-bounds-in-the-original-value-coordinates}

Keep \texttt{S=c+iy}, \texttt{c=k/2}, and the actual mass \nolinkurl{M_sigma=sqrt(2pi)} of \nolinkurl{sigma(y)=|Gamma(1/4+iy/2)|^2/(2pi)}. In root-value coordinates put \nolinkurl{y_alpha=(beta_alpha-c)/i}, so \nolinkurl{|y_alpha|<=R_y=k sqrt(delta^2+gamma^2)}. At the four cutoffs \texttt{q-1,q,2q-1,2q} let C\_N be the full Gamma covariance with degrees zero through N and let \nolinkurl{G_N^val=C_N^{-1}}.

The Cauchy estimate from the exact Meixner--Pollaczek generating function gives, for N\textgreater=1 and j\textless=N, \[
|p_j(y_\alpha)|^2/\gamma_j
\le\frac{e^2}{M_\sigma}(N+1)^{1/2}(2N+1)^{R_y+1}.
\] Here are the finite details. In \nolinkurl{sum_j p_j(y)t^j/j!=(1+t^2)^{-1/4} exp(y arctan t)} take \texttt{r=N/(N+1)}. The convergent arctangent series gives \nolinkurl{|arctan t|<=arctanh r=(1/2)log(2N+1)} on that circle, and \texttt{1-r\^{}2\textgreater{}=1/(N+1)}. Cauchy's coefficient integral, together with \texttt{r\^{}\{-j\}\textless{}=e}, bounds \texttt{\textbar{}p\_j(y)\textbar{}/j!} by \nolinkurl{e(N+1)^{1/4}(2N+1)^{R_y/2}}. Finally \nolinkurl{j!/(1/2)_j<=2j+1<=2N+1}: its induction starts at one, and multiplication by \texttt{(j+1)/(j+1/2)} takes \texttt{2j+1} to \texttt{2j+2}. Squaring the coefficient estimate and inserting \nolinkurl{gamma_j=M_sigma j!(1/2)_j} proves the displayed bound. It follows by summing the norms of every value column that \[
\lambda_{\max}(C_N)\le
A_N^{\rm val}:=\frac{qe^2}{M_\sigma}(N+1)^{3/2}
 (2N+1)^{R_y+1}.
\tag{KF33}
\] For a lower bound retain the first q columns. The Lagrange polynomial at a root alpha has numerator \nolinkurl{prod_{beta!=alpha}(y-y_beta)} and denominator \nolinkurl{prod_{beta!=alpha}(y_alpha-y_beta)}. Every root separation has modulus at least \texttt{2delta}. Repeated use of the exact Gamma bound \nolinkurl{||yP||<=2(deg P+1)||P||} gives \[
\|\ell_\alpha\|_\sigma
\le\sqrt{M_\sigma}
\left(\frac{2q+R_y}{2\delta}\right)^{q-1}.
\] To check the multiplier constant, in the orthonormal Gamma frame the two off-diagonal coefficients of y are \nolinkurl{sqrt((j+1)(j+1/2))} and \texttt{sqrt(j(j-1/2))}. The raising part on degrees through n has norm at most n+1 and the lowering part has norm at most n.~Their sum is at most \texttt{2(n+1)}. Multiplication by each factor \texttt{y-y\_beta} therefore has norm at most \texttt{2q+R\_y} during the q-1 factors in a Lagrange numerator. This proves the asserted iteration starting from the constant polynomial of norm \texttt{sqrt(M\_sigma)}. Thus interpolation from a Euclidean unit root-value vector has squared norm at most \[
B_k^{\rm val}:=qM_\sigma
 \left(\frac{2q+R_y}{2\delta}\right)^{2(q-1)}.
\tag{KF34}
\] The minimum through degree N has no larger norm than this interpolation trial, so \nolinkurl{G_N^val<=B_k^val I}. KF33 gives its lower bound. With \[
\kappa_k^*=\max\{1,A_{2q}^{\rm val}B_k^{\rm val}\},
\] all four complete Gamma metrics satisfy \[
(A_{2q}^{\rm val})^{-1}I\preceq G_N^{\rm val}
\preceq B_k^{\rm val}I,
\quad \operatorname{cond}(G_N^{\rm val})\le\kappa_k^*,
\quad\log\kappa_k^*=O_h(q\log(q+2)).
\tag{KF35}
\] The estimate is on the full covariance before restricting or quotienting. No original value direction is removed. The phases relating S and y have modulus one in each polynomial degree and root values are unchanged.

By comparison of distances under two positive metrics, KF31 implies the actual Gamma angle bound \[
\operatorname{dist}_{G_N}(l,K)
\ge\frac{\varepsilon_k}{\sqrt{\kappa_k^*}}
 \operatorname{dist}_{G_N}(l,K\cap\mathcal L),\qquad l\in\mathcal L.
\tag{KF36}
\] Indeed the first metric distance is at least the square root of its smallest eigenvalue times Euclidean distance. Use KF31 and then \nolinkurl{dist_E(l,K_0)>=dist_G(l,K_0)/sqrt(lambda_max(G))}. This proves the stated ratio.

\subsection{6. Removing the fixed low graph from the actual kernel determinant}\label{removing-the-fixed-low-graph-from-the-actual-kernel-determinant}

Put \nolinkurl{K_0=K intersect L}, \texttt{s\_k=dim\ K\_0}, and let K\_bar be the image of K in V\_A/L. Its rank is \texttt{h=m-s\_k}. All spaces and frames in this paragraph are fixed across the four cutoffs. First give K\_bar the quotient metric from K/K\_0, denoted H\_bar\^{}K. Second give it the restriction of the complete quotient metric from V\_A/L, denoted H\_bar\^{}V. Their definitions minimize over K\_0 and L respectively.

There is the exact comparison \[
0\le e_N^{\rm ang}:=
\log\det H_{\rm bar,N}^K-\log\det H_{\rm bar,N}^V
\le v\log\kappa_k^*+2v\log(1/\varepsilon_k)=:B_k^{\rm ang}.
\tag{KF37}
\] Here is the complete subspace argument. In the original Gamma metric remove K\_0 orthogonally, put \nolinkurl{A=K orthogonal-minus K_0}, \nolinkurl{B=L orthogonal-minus K_0}, and put \nolinkurl{theta=epsilon_k/sqrt(kappa_k^*)}. The map from K/K\_0 to V\_A/L is \texttt{J=(I-P\_B)\textbar{}\_A}. KF36, applied to each b in B, gives \[
\|(I-P_A)b\|^2\ge\theta^2\|b\|^2,
\quad\|P_A|_B\|^2\le1-\theta^2,
\quad(P_B|_A)^*=P_A|_B.
\] Consequently, with \texttt{C=P\_B\textbar{}\_A}, \[
J^*J=I_A-C^*C\succeq\theta^2 I_A,
\qquad \operatorname{rank}(C^*C)\le\dim B\le v.
\] There are thus at most v eigenvalues different from one, and every eigenvalue belongs to \texttt{{[}theta\^{}2,1{]}}. In a frame orthonormal for H\_bar\^{}K the second metric is exactly J\^{}*J. It follows that \[
0\le e_N^{\rm ang}=-\log\det(J^*J)
\le-2v\log\theta
=v\log\kappa_k^*+2v\log(1/\varepsilon_k).
\] This proves KF37, including nontrivial intersection K\_0. The fixed rank is used only after the actual quantitative angle bound has been proved.

Choose a fixed basis of K\_0 and fixed lifts completing it to a frame of K. Schur complementation of its complete restricted Gram gives \[
\log\det H_{K,N}=\log\det H_{K_0,N}
 +\log\det H_{\rm bar,N}^K+C_{\rm frame},
\tag{KF38}
\] where the change from the original I\_K frame has its fixed determinant recorded by C\_frame. It cancels in the four-return only because that same frame is used at all four cutoffs.

The complete source minima decrease with N. Their restriction to K\_0 therefore decreases, while KF35 bounds all four by the same two scalar factors in a Euclidean-orthonormal frame. Pairing the signs gives \[
0\le\mathcal R\log\det H_{K_0,N}
\le2s_k\log\kappa_k^*.
\tag{KF39}
\] Also KF37 gives \nolinkurl{|R e_N^ang|<=2 B_k^ang}, because each of the four errors belongs to the same interval \texttt{{[}0,B\_k\^{}ang{]}}. Equations KF38--KF39 therefore prove \[
\left|\mathcal R\log\det H_{K,N}^{\sigma}
-\mathcal R\log\det H_{\rm bar,N}^{V,\sigma}\right|
\le2s_k\log\kappa_k^*+2B_k^{\rm ang}
=O_{\rm actual}(q\log(q+2))=o(kq).
\tag{KF40}
\] This removes the finite low-degree component with a proved error for the actual invariant kernel. It is stronger than deleting a rank-v term without controlling its angle.

\subsection{7. The complete relation graph and its source transport}\label{the-complete-relation-graph-and-its-source-transport}

For all four original cutoffs put \texttt{L=N-q}; it takes values \texttt{-1,0,q-1,q}. Retain the exact outer products for every original shift, \[
d_{rs}(S')=\prod_{\substack{0\le i,j\le k\\
(i,j)\notin\{r,\ldots,r+k'\}\times\{s,\ldots,s+k'\}}}
(S'+b_{rs}-\beta_{ij}),
\] and the whole relation operator \[
\mathcal U_{A,L}h(S')=\sum_{r,s}a_{rs}d_{rs}(S')h(S'+b_{rs}),
\quad \mathcal T_A(\chi h)=\chi'\mathcal U_{A,L}h.
\tag{KF41}
\] It is injective and has degree \texttt{deg\ h+g}, with leading coefficient \nolinkurl{mu_v binom(q+deg h,v)} times the leading coefficient of h. Thus \[
Y_N=\mathcal P'_{L+g}/\mathcal U_{A,L}\mathcal P_L
\tag{KF42}
\] has dimension g at every cutoff. Negative degree means zero, so at the first cutoff it is exactly P'\_\{g-1\}. Its source norm is the complete weighted lower norm \[
\|\eta\|_{\chi',1}^2=
\int_{\mathbb R}|\chi'(c'+iy)\eta(c'+iy)|^2\sigma(y)\,dy,
\qquad c'=k/2-4.
\tag{KF43}
\] The source order is still one, and its mass factor has not been changed.

The actual map is \[
\Psi_N:V_A/\mathcal L\longrightarrow Y_N,
\qquad [p]\longmapsto[\mathcal T_Ap/\chi'].
\tag{KF44}
\] If a lift changes by chi h, KF41 changes its image by U\_Ah. If its image is zero, subtract that same chi h and KF4 leaves a polynomial in L; this proves injectivity. Surjectivity follows by solving the onto T\_A source equation. The unique degree-below-q representative gives eta of degree below g by KF7, so the coordinate matrix of Psi\_N on these low target representatives is independent of N. The exact actual plane in Y\_N is \[
\overline K_N=\Psi_N(K/ K_0).
\tag{KF45}
\] Its fixed low-polynomial frame is obtained from I\_K by KF7 and quotient by its actual nullspace K\_0; equivalently it is the projection of the kernel of the full residual constraint KF10. Thus KF45 retains the complete invariant image rather than an arbitrarily oriented rank-h plane in Y\_N.

Here are finite constants for the source transport used in this map. Set D=N-v, let B\_D\^{}+ be the original WCF12 order-one bound \[
B_D^+=e2^v A_{\rm abs}e^{2\pi\gamma}(2D+1)^{4\delta}
\sqrt{2(D+v)+1}\quad(D\ge1),\quad M_D=\max(1,B_D^+),
\] and put \[
g_0=(-i)^v\mu_v/v!,\quad J_D=(D+1)\log^+(1/|g_0|),\quad
E_N^*= (N+1)\log M_D+J_D.
\tag{KF46}
\] All present endpoints have D\textgreater=1. The exact orthonormal source operator \nolinkurl{S_Np=(P_<v p,T_Ap)} has a block I\_v and the WCF upper triangular block. Its forward norm is at most M\_D. Its determinant is exactly \[
\det A_{D,1}=g_0^{D+1}\prod_{i=0}^D\rho_{i+v}/\rho_i,
\qquad \rho_i=\sqrt{i!/(1/2)_i}.
\] Since rho\_i increases, \nolinkurl{|det S_N|^{-1}<=exp J_D}. The singular-value complement identity therefore proves, for every rank h, \[
\log\|\bigwedge^h S_N\|\le h\log M_D,\qquad
\log\|\bigwedge^hS_N^{-1}\|\le(N+1-h)\log M_D+J_D.
\tag{KF47}
\] These are bounds for the full original source map, not scalar bounds multiplied by a changing relation rank.

Its exact image of the conductor preimage is \nolinkurl{C^v direct-sum chi' P'_{L+g}}. Its image of the original relation chi h is \nolinkurl{(P_<v(chi h),chi' U_Ah)}. After quotient by the entire L, the free first component vanishes and the remaining metric is exactly the weighted complete quotient KF42--KF43. All low graph entries were retained before taking this quotient. The induced map on any actual subspace, in particular KF45, is a minimal lift followed by S\_N and an orthogonal quotient projection. Projection contracts every exterior norm, and the same argument for S\_N\^{}\{-1\} supplies the reverse bound. Thus, in the transported actual frame, \[
\left|\log\det H_{\overline K_N,Y_N}
-\log\det H_{\rm bar,N}^{V,\sigma}\right|\le2E_N^*.
\tag{KF48}
\] This proves the source-graph comparison for the actual plane, using the complete original graph and its quotient order. WCF gives \nolinkurl{E_N^*=O_actual(q log(q+2))=o(kq)} at all four cutoffs.

Combining KF40 and KF48 with the previously proved JM25/MR21 arithmetic comparison gives the finite final receiver \[
\boxed{\begin{aligned}
\left|\mathcal R\log\det H_{K,N}^{\rm ar}
-\mathcal R\log\det H_{\overline K_N,Y_N}\right|
\le{}&2mL_k^{\rm form}+2s_k\log\kappa_k^*
+2B_k^{\rm ang}+2\sum_N E_N^*\\
={}&o_{\rm actual}(kq).
\end{aligned}}
\tag{KF49}
\] The sum is over exactly \texttt{q-1,q,2q-1,2q}. Their source relation degrees are exactly \texttt{-1,0,q-1,q}, and their target polynomial degrees are \texttt{g-1,g,q+g-1,q+g}. The arithmetic source retains \nolinkurl{w_h=|2xi/h|^2/(2pi)} and its actual complete divisor. Its physical units enter I\_K and the same original maps throughout; their prior isometries, rather than a replacement of the units by constants, are used by JM25.

\subsection{8. What has been evaluated and the remaining connecting map}\label{what-has-been-evaluated-and-the-remaining-connecting-map}

KF11--KF20 prove exact high-degree transversality modulo the actual low intersection and explicit growth bounds for the sum of the actual projection minors. KF25--KF32 evaluate that low intersection through a fixed finite jet matrix on each residue class and prove its quantitative nonzero-angle bound. KF49 proves that the entire fixed low component and the conductor source transport have subleading four-return cost.

The remaining rank \texttt{m-s\_k} metric is the restriction in KF45 of the complete quotient KF42. The relations are the actual sum of shifted outer products KF41. Its exact multiplication defect is \[
\mathcal U_A(Sh)-S'\mathcal U_Ah
=\sum_{r,s}a_{rs}b_{rs}d_{rs}(S')h(S'+b_{rs}),
\tag{KF50}
\] with all original coefficients present. The identity alone does not prove that its right side has nonzero class modulo the complete relation image. Thus closure under multiplication is neither assumed nor inferred from this display. The following exact graph comparison decides the relation to an outer polynomial ideal without making that inference. It is not legitimate to replace U\_A by one d\_rs on the grounds that their source and target dimensions agree. Nor do KF18--KF20 convert coefficient control into a Gamma minimum without the metric steps above.

At each cutoff the residual determinant can be evaluated without dropping any cross Gram as follows. Choose the fixed low eta frame B\_K for KF45; let H\_N be the full weighted monomial Gram on P'\_\{L+g\}, let U\_N be the full coefficient matrix of KF41, and let E\_N insert B\_K in that source. Then its exact restricted minimum is \[
\boxed{H_{\overline K_N,Y_N}=
E_N^*H_NE_N-E_N^*H_NU_N(U_N^*H_NU_N)^{-1}U_N^*H_NE_N.}
\tag{KF51}
\] At L=-1 the second term is zero. The inverted matrix is positive definite because U\_A is injective. The formula follows by minimizing the full quadratic form of \texttt{E\_Nb+U\_Nh} in h; the minimizing coefficient is \nolinkurl{-(U_N^*H_NU_N)^{-1}U_N^*H_NE_Nb}. It retains the complete source-graph correction at every endpoint. KF49 shows that evaluating its four-return determines the original kernel's kq coefficient, with the proved error, but the present argument does not yet give that remaining coefficient.

\subsection{9. An exact graph over the actual outer polynomial}\label{an-exact-graph-over-the-actual-outer-polynomial}

There is a canonical outer polynomial for this comparison, retaining every shift: put \texttt{d\_A=U\_A1}. By KF41 its degree is exactly g and its leading coefficient is \nolinkurl{a_A=mu_v binom(q,v)}, which is nonzero. For each source degree L\textgreater=0 perform Euclidean polynomial division, with that actual leading coefficient, on the whole operator U\_A: \[
\mathcal U_Ah=d_A\,\mathcal C_Lh+\mathcal R_Lh,
\qquad \deg\mathcal R_Lh<g,\quad h\in\mathcal P_L.
\tag{KF52}
\] Here C\_L and R\_L are the unique linear maps given by division. They retain all original coefficients in KF41. If h has degree n, the leading term calculation in KF41 proves that C\_L h has degree n and leading coefficient multiplied by \nolinkurl{binom(q+n,v)/binom(q,v)}. Therefore C\_L is upper triangular in the increasing monomial bases and \[
\det\mathcal C_L=
\prod_{n=0}^L\frac{\binom{q+n}{v}}{\binom qv}>0,
\qquad
0\le\log\det\mathcal C_L
\le v(L+1)\log\frac{q+L}{q-v+1}.
\tag{KF53}
\] For v=0 the determinant is one and the logarithmic upper bound is zero. For v\textgreater=1 each ratio is the product, over \texttt{0\textless{}=j\textless{}v}, of \texttt{(q+n-j)/(q-j)}, between one and \nolinkurl{((q+L)/(q-v+1))^v}. This proves KF53 without estimating entries by a rank-only replacement. At L\textless=q the upper bound is \texttt{O\_actual(q)}.

Define the complete graph coefficient map \nolinkurl{F_L=R_L C_L^{-1}:P_L->P'_{g-1}}. Equation KF52 proves exactly \[
\mathcal U_A\mathcal P_L
=\{d_Aw+F_Lw:w\in\mathcal P_L\}.
\tag{KF54}
\] Every polynomial of degree at most L+g has a unique decomposition \texttt{r+d\_Aw} with \texttt{deg\ r\textless{}g}, \texttt{deg\ w\textless{}=L}. The linear isomorphism \[
\mathcal J_L(r+d_Aw)=r-F_Lw+d_Aw,
\quad
\mathcal J_L^{-1}(r+d_Aw)=r+F_Lw+d_Aw
\tag{KF55}
\] fixes each low representative r and sends the complete relation image in KF54 onto the entire multiplication image \texttt{d\_A\ P\_L}. Its determinant in this displayed decomposition is one; its off-diagonal block is the full F\_L. This is the exact relation between the original conductor quotient and that outer polynomial quotient. In particular the actual low frame B\_K is preserved, and all of its cross terms remain.

To carry the metric as well, equip the outer polynomial source with the Gram \nolinkurl{(J_L^{-1})^* H_N J_L^{-1}}. Then J\_L is an isometry from the complete weighted source H\_N to that displayed source: direct substitution gives equality of squared norms. It descends to an isometry from Y\_N onto the quotient by \texttt{d\_A\ P\_L} with this exact transformed metric. Replacing this transformed Gram by H\_N would discard F\_L and is not a consequence of determinant one. At L=-1 set J\_L to the identity, so this isometry includes the first cutoff.

Equations KF52--KF55 are finite evaluations of the complete source-graph correction as polynomial division and an invertible shear. In particular the exact obstruction to deleting that correction is the metric action of F\_L, not the scalar determinant in KF53. The present proof establishes no subleading bound for that metric action. KF51 and KF55 give two equivalent exact next calculations, with all four cutoffs and the full actual frame, while KF49 has already bounded the low-angle and conductor transport corrections by \texttt{o(kq)}.

\section*{Exact programme proof sources}
\par\href{https://github.com/KokunoYumeto/zeta-function-research-reader/blob/5992ad50c0f2a06921f1fcaded7b4e38097c0739/workbenches/splitzero-tandem/continuations/20260920-original-kernel-mixed-determinants/ORIGINAL_KERNEL_MIXED_DETERMINANT.tex#L36}{ORIGINAL\_KERNEL\_MIXED\_DETERMINANT.tex, MR1 and the following complete proof.}
\par\href{https://github.com/KokunoYumeto/zeta-function-research-reader/blob/5992ad50c0f2a06921f1fcaded7b4e38097c0739/workbenches/splitzero-tandem/continuations/20260920-original-kernel-mixed-determinants/providers/ORIGINAL_CONDUCTOR_STRIP_FRAME.tex#L27}{ORIGINAL\_CONDUCTOR\_STRIP\_FRAME.tex, OCF1 and the following complete proof.}
\par\href{https://github.com/KokunoYumeto/zeta-function-research-reader/blob/5992ad50c0f2a06921f1fcaded7b4e38097c0739/workbenches/splitzero-tandem/continuations/20260920-original-kernel-mixed-determinants/providers/WEIGHTED_CONDUCTOR_FORWARD.tex#L28}{WEIGHTED\_CONDUCTOR\_FORWARD.tex, WCF1 and the following complete proof.}
\par\href{https://github.com/KokunoYumeto/zeta-function-research-reader/blob/5992ad50c0f2a06921f1fcaded7b4e38097c0739/workbenches/splitzero-tandem/continuations/20260920-original-kernel-mixed-determinants/providers/09_INPUT_CUMULATIVE_PROOFS.tex#L19060}{09\_INPUT\_CUMULATIVE\_PROOFS.tex, CEP2 and the following complete proof.}
